\section{Related Work}
\label{sec:related_work}

\subsection{Cognitive Models of Second Language Acquisition}
SLA research provides theoretical scaffolding for evaluating neural networks. Krashen’s Input Hypothesis \cite{krashen1985input} parallels the unsupervised pre-training of LLMs, where acquisition relies on comprehensible input. Connectionist theories, such as MacWhinney’s Competition Model \cite{macwhinney2001competition} and Ellis’s work on emergentism \cite{ellis2002frequency}, argue that language learning is a statistical, frequency-driven process. These frameworks help analyze whether Transformers exhibit human-like transitional phases—like vocabulary mixing—when acquiring a new language.

\subsection{LLMs as Cognitive and Linguistic Models}
Large Language Models increasingly serve as subjects for psycholinguistic study. Linzen and Baroni \cite{linzen2021syntactic} demonstrate that deep neural networks capture emergent syntax without explicit rule programming. While cross-lingual transfer is well-documented \cite{conneau2019cross}, fine-tuning monolingual models on new languages often triggers \textit{catastrophic forgetting}. Mitigation techniques, such as those formalized by Kirkpatrick et al. \cite{kirkpatrick2017overcoming}, are central to our methodology for retaining native English capabilities while integrating Spanish.

\subsection{Gamification and Interactive Educational Tools}
Effective educational tools rely on engagement and immediate feedback \cite{plass2015foundations}. Interactive virtual environments lower the learner's ``affective filter'' \cite{godwin2014games}, promoting exploration. By utilizing the Godot engine for a physics-based simulation—where linguistic inaccuracies physically add weight to a virtual boat—our tool leverages dynamic feedback and embodied cognition \cite{kapp2012gamification}. This allows researchers to visually and interactively assess the evolving capability of the model in real-time.